\documentclass[11pt,a4paper]{article}
\usepackage[margin=25mm]{geometry}
\usepackage{amsmath,amssymb,amsthm,mathtools,bm}
\usepackage{newtxtext,newtxmath}
\usepackage{booktabs,array,microtype}
\usepackage[hidelinks]{hyperref}
\usepackage{enumitem}
\setlist{nosep}
\setlength{\parskip}{4pt}
\setlength{\parindent}{1.2em}
\allowdisplaybreaks[1]
\newtheorem{proposition}{Proposition}
\DeclareMathOperator{\ran}{ran}
\DeclareMathOperator{\rank}{rank}
\DeclareMathOperator{\supp}{supp}
\newcommand{\Tr}{^{\mathsf T}}
\newcommand{\ip}[2]{\left\langle #1,#2\right\rangle_{\mathrm F}}
\newcommand{\Wa}{\mathcal W}
\newcommand{\Ta}{\mathcal T}
\newcommand{\Ha}{\mathcal H}
\newcommand{\Ea}{\mathcal E}
\newcommand{\Id}{\mathcal I}
\renewcommand{\vv}[1]{\mathbf{#1}}
\title{Exact Orthogonal-CI Contraction for Arbitrary Selected\newline
Valence-Bond Structure Spaces}
\author{Theory and implementation note}
\date{16 October 2026}

\begin{document}
\maketitle
\begin{abstract}
An arbitrary selected valence-bond (VB) structure space is generally not
invariant under orbital orthogonalization. This does not prevent exact
evaluation of its Hamiltonian and overlap actions through an orthogonal
configuration-interaction (CI) representation. The transformed wave function
must be retained, but its complete spin-string product array need not be
formed. We derive an exact source-factor representation, apply the
Slater--Condon one-body generators separately to the two spin factors, and
contract only the original products needed for the final structure-space
projection. The resulting equations preserve the selected variational space
without a closure assumption or intermediate amplitude truncation. Source
and channel tiling bound the numerical workspace. Optional spin-exchange
reduction and causal pruning of the exterior adjoint are derived separately.
The scope is the Hamiltonian/overlap action, its active-integral adjoint,
and the factor representation of the directional integral adjoint.
\end{abstract}

\section{The selected VB space}

Let $n$ be the number of active spatial orbitals, with $n_\alpha$ and
$n_\beta$ active electrons of each spin. The complete one-spin string
spaces have dimensions
\begin{equation}
 U_\alpha=\binom{n}{n_\alpha},\qquad
 U_\beta=\binom{n}{n_\beta}.
\end{equation}
An index $a$ identifies an $\alpha$ string and $b$ a $\beta$ string.
For a selected VB structure $I$, let $w_{ab,I}$ be its signed expansion
coefficient on the corresponding spin-string product, with all normalization
conventions included. For structure coefficients $\vv C$, define
\begin{equation}\label{eq:W}
 X=\Wa\vv C,\qquad X_{ab}=\sum_I w_{ab,I}C_I.
\end{equation}
Here $\Wa$ is a sparse linear map from structure coefficients to a matrix
of spin-string product coefficients. The matrix notation does not prescribe
a dense allocation: the map can be stored as sparse incidences between
structures and touched string products. It introduces no determinant-pair
traversal or separately optimized determinant coefficient vector.

The original support and its two marginal index sets are
\begin{align}
 \Omega&=\{(a,b):w_{ab,I}\ne0\text{ for at least one selected }I\},
 &t&=|\Omega|,\\
 A_\Omega&=\{a:\exists b,\ (a,b)\in\Omega\},&
 J&=\{b:\exists a,\ (a,b)\in\Omega\},\qquad r=|J|.
\end{align}
Support is determined by the selected structures, not by numerical
smallness or cancellations in a particular Davidson vector.
Using the Frobenius inner product, the adjoint structure map is
\begin{equation}\label{eq:Wadj}
 (\Wa^*Z)_I=\sum_{(a,b)\in\Omega}w_{ab,I}Z_{ab}.
\end{equation}
Thus only entries in $\Omega$ are needed at the \emph{end} of a complete
action. Nothing in this statement permits dropping intermediate entries
outside $\Omega$.

\section{Orthogonalization without a closure assumption}

At a fixed orbital point, write the nonorthogonal active orbitals as
$\phi=\chi R$, where $\chi$ is an orthonormal basis of the same active
one-particle space. For example, if $O$ is their positive-definite overlap,
\begin{equation}
 O=R\Tr R,\qquad \chi=\phi R^{-1}.
\end{equation}
This is an evaluation-coordinate change, not a change to the physical
localization constraints or to the orbital optimization variables.
Multilinearity of antisymmetrized products gives the exterior maps
\begin{equation}
 T_\sigma=\wedge^{n_\sigma}R,\qquad
 (T_\sigma)_{a i}=\det R[a,i],\qquad \sigma\in\{\alpha,\beta\},
\end{equation}
where the row and column subsets contain $n_\sigma$ orbitals with a fixed
ordering convention. The orthogonal-basis coefficient matrix is
\begin{equation}\label{eq:T}
 Y=\Ta X=T_\alpha X T_\beta\Tr,\qquad
 \Ta^* Z=T_\alpha\Tr Z T_\beta.
\end{equation}
The minor matrices $T_\sigma$ are mathematical definitions; the
implementation applies elementary exterior transforms without assembling
them densely.

Let $\Ha$ be the active Hamiltonian action in the orthogonal spin-string
product basis. The selected structure Hamiltonian and overlap are exactly
\begin{equation}\label{eq:restriction}
 \boxed{\quad H=\Wa^*\Ta^*\Ha\Ta\Wa,\qquad
 M=\Wa^*\Ta^*\Ta\Wa.\quad}
\end{equation}
An orbital-dependent scalar inactive-plus-nuclear energy $E_0$ contributes
$E_0M$ to the total structure Hamiltonian and is separate from the active
action $H$ in this note.
The variational energy is
\begin{equation}
 E=E_0+\frac{\vv C\Tr H\vv C}{\vv C\Tr M\vv C},
\end{equation}
with $H$ denoting the active contribution of Eq.~\eqref{eq:restriction}.

\begin{proposition}[Exact selected-space restriction]\label{prop:closure}
For any selected structure map $\Wa$, Eq.~\eqref{eq:restriction} gives
the exact selected-space Hamiltonian and overlap, whether or not
$\Ta\ran\Wa\subseteq\ran\Wa$.
\end{proposition}
\begin{proof}
For arbitrary structure vectors $\vv u$ and $\vv v$, the two physical
states have orthogonal-basis amplitudes $\Ta\Wa\vv u$ and
$\Ta\Wa\vv v$. Their Hamiltonian matrix element is
\begin{equation}
 \ip{\Ta\Wa\vv u}{\Ha(\Ta\Wa\vv v)}
 =\vv u\Tr\bigl(\Wa^*\Ta^*\Ha\Ta\Wa\bigr)\vv v.
\end{equation}
Replacing $\Ha$ by the identity gives the overlap. No intermediate
restriction to $\ran\Wa$ was inserted.
\end{proof}

The nonclosure issue arises only if a different operation is used, for
example by zeroing transformed amplitudes outside the original selected
support before applying $\Ha$. That operation changes the wave function.
Equation~\eqref{eq:restriction} does not do so. Davidson vectors and Ritz
coefficients remain in the selected VB structure space; the orthogonal
representation is used internally to evaluate their images.
For linearly independent selected structures, their eigenproblem remains
\begin{equation}
 H\vv C=E_{\rm act}M\vv C.
\end{equation}
The adjoint $\Wa^*$ contracts against the selected test structures; it is
not an inverse coefficient transform. The nonidentity overlap $M$ is
retained, so this procedure does not replace the selected VB problem by
an unconstrained orthogonal-CI eigenproblem.

\subsection{A minimal nonclosed example}

Consider a two-by-two coefficient matrix with the single selected product
$X=e_2e_2\Tr$, and let
\begin{equation}
 T_\alpha=\begin{pmatrix}1&a\\0&1\end{pmatrix},\qquad
 T_\beta=\begin{pmatrix}1&b\\0&1\end{pmatrix}.
\end{equation}
This is a linear-algebra example of the coefficient map, not an assertion
that a single spin product is a spin-adapted VB structure. Its transformed
amplitudes are
\begin{equation}\label{eq:example}
 Y=\begin{pmatrix}ab&a\\b&1\end{pmatrix}
 =\begin{pmatrix}a\\1\end{pmatrix}
  \begin{pmatrix}b\\1\end{pmatrix}\Tr.
\end{equation}
For nonzero $a$ and $b$, all four products are present, so the original
one-dimensional selected space is not closed. Nevertheless, two vectors
represent the complete $Y$ exactly. The overlap is
$(1+a^2)(1+b^2)$, not the value $1$ obtained by retaining only $Y_{22}$.
An exact Hamiltonian contraction must likewise use all transformed
amplitudes, explicitly or implicitly. The factorized route below retains
them implicitly.

\section{Exact source factors}

\subsection{Column decomposition}

Order the original touched $\beta$ columns as $J=(j_1,\ldots,j_r)$, and set
\begin{equation}\label{eq:source}
 A_{:\ell}=X_{:j_\ell},\qquad B_{:\ell}=e_{j_\ell},\qquad
 X=AB\Tr.
\end{equation}
This is an identity for every input vector. It is not a truncated singular
value decomposition, and $r$ is not a fitted rank. Transforming only the
two factors gives
\begin{equation}\label{eq:factors}
 L=T_\alpha A,\qquad F=T_\beta B,\qquad
 \Ta X=LF\Tr.
\end{equation}
The factors have dimensions $U_\alpha\times r$ and $U_\beta\times r$.
The complete transformed $U_\alpha\times U_\beta$ matrix need not be
formed, even if every one of its entries is nonzero.

The distinction between $r$ and the number of structures matters. A single
covalent structure with $m$ disjoint singlet pairs generally touches
$2^m$ original source columns. Several structures can share most of those
columns, and a broad selected space can touch all $U_\beta$ columns.
Thus $\rank X\le r$ is an upper bound, not an identification of $r$ with
one, with $\rank X$, or with the number of structures.

\subsection{A minimum mixed row/column cover}\label{sec:cover}

Grouping every source on the $\beta$ side is not always economical. View
the original input support as a bipartite graph: vertices are the touched
$\alpha$ and $\beta$ strings, and each supported product $(a,b)$ is one
edge. Choose a vertex cover $C_\alpha\cup C_\beta$, so that every edge
touches at least one chosen vertex. Assign an edge to its $\alpha$ endpoint
if that endpoint is in $C_\alpha$; otherwise assign it to its covered
$\beta$ endpoint. With this disjoint assignment,
\begin{equation}\label{eq:cover}
 X=\sum_{a\in C_\alpha}e_a u_a\Tr
   +\sum_{b\in C_\beta}v_b e_b\Tr=AB\Tr.
\end{equation}
Here $(u_a)_b=X_{ab}$ only for edges assigned to row $a$, and
$(v_b)_a=X_{ab}$ only for edges assigned to column $b$; the remaining
entries are zero. Thus no coefficient is counted twice, including edges
whose two endpoints are both covered. Each chosen vertex defines one
coordinate-anchored outer product.

Let $\nu(\Omega)$ be the maximum matching size, i.e., the largest number
of edges with no shared endpoints. A minimum cover has cardinality
\begin{equation}
 r_* = |C_\alpha|+|C_\beta|=\nu(\Omega)
 \le\min(|A_\Omega|,|J|).
\end{equation}
For completeness, the cover can be certified directly from a maximum
matching. Start at unmatched $\alpha$ vertices and follow alternating
unmatched/matched edges. If the reached vertex sets are $Z_\alpha$ and
$Z_\beta$, then
$(A_\Omega\setminus Z_\alpha)\cup Z_\beta$ covers every edge and
contains exactly one endpoint per matched edge. Any cover must contain
at least one endpoint of each disjoint matched edge, establishing
minimality. The implementation uses layered augmenting paths, with
$O(t\sqrt{U_\alpha+U_\beta})$ worst-case planning work and linear graph
storage. The plan is independent of the numerical Davidson vector.

This is a minimal factor count for \emph{support-only coordinate-anchored}
decompositions, not a numerical rank claim. In particular,
$\rank X\le\nu(\Omega)$; additional coefficient relations can lower the
actual rank further. A diagonal support needs one source per edge. In
contrast, the support
\begin{equation}
 \{(a_0,b_j):1\le j\le m\}
 \cup\{(a_i,b_0):1\le i\le m\}
\end{equation}
needs only two mixed sources, although either a pure row or pure column
decomposition needs $m+1$ sources. This reduces the factor dimension,
rather than enlarging a tile or imposing an amplitude threshold.

All subsequent Hamiltonian and sampling equations apply unchanged, with
$r=r_*$ and the factors of Eq.~\eqref{eq:cover}. When a common exchange
parity is used, construct the graph from the \emph{representative input}
support of Section~6.1; the final projection still uses the complete
original $\Omega$. If the original column cover is already minimal, its
ordering may be retained to avoid changing established summation order.

\section{Hamiltonian action on the factors}

\subsection{One-body channels and the complete active Hamiltonian}

Use real spatial orbitals and real, spin-independent molecular integrals.
Let $e_{pq}=\sum_\sigma a^\dagger_{p\sigma}a_{q\sigma}$, and use the
chemists' notation $(pq|rs)$ for two-electron integrals. The active
Hamiltonian is
\begin{equation}\label{eq:Hsecond}
 \hat H=\sum_{pq}h_{pq}e_{pq}
 +\frac12\sum_{pqrs}(pq|rs)
       \bigl(e_{pq}e_{rs}-\delta_{qr}e_{ps}\bigr).
\end{equation}
Here $h$ includes the effective one-electron contribution of the inactive
orbitals. Define
\begin{equation}
 f_{ps}=h_{ps}-\frac12\sum_q(pq|qs),\qquad
 N=n_\alpha+n_\beta>0.
\end{equation}
Packed pair labels $P=(p,q)$ have $p\ge q$, so there are
$d=n(n+1)/2$ channels. Define the symmetric generators
\begin{equation}
 \hat E_{\sigma P}=\begin{cases}
 a^\dagger_{p\sigma}a_{p\sigma},&p=q,\\
 a^\dagger_{p\sigma}a_{q\sigma}
 +a^\dagger_{q\sigma}a_{p\sigma},&p>q,
 \end{cases}
 \qquad \hat E_P=\hat E_{\alpha P}+\hat E_{\beta P}.
\end{equation}
Let $f_P=f_{pq}$, $V_{PQ}=(pq|rs)$ for $Q=(r,s)$, and $d_P=1$ for
diagonal pairs and zero otherwise. A symmetric kernel containing the
\emph{entire} active Hamiltonian is
\begin{equation}\label{eq:K}
 K_{PQ}=\frac12 V_{PQ}
       +\frac{f_Pd_Q+d_Pf_Q}{2N},\qquad
 \hat H=\sum_{PQ}K_{PQ}\hat E_P\hat E_Q.
\end{equation}
Indeed, $\sum_Pd_P\hat E_P=\hat N=N\hat I$ on the fixed-electron
space. The second term in $K$ therefore gives
$\sum_P f_P\hat E_P$; the first term supplies the quadratic part of
Eq.~\eqref{eq:Hsecond}. The usual integral permutation symmetries account
for both ordered off-diagonal orbital pairs through $\hat E_P$; no extra
off-diagonal multiplicity is needed in $K$.

Let $E_{\sigma P}$ be the matrix of $\hat E_{\sigma P}$ on its
one-spin string space. The Slater--Condon rules make this operation a
sparse signed occupation/replacement scatter, rather than a dense string
pair contraction. Its action on a product matrix is
\begin{equation}
 \Ea_P(Z)=E_{\alpha P}Z+Z E_{\beta P}\Tr,\qquad
 \Ha=\sum_{PQ}K_{PQ}\Ea_P\Ea_Q.
\end{equation}
These identities require a complete carrier for each fixed-spin sector,
or an exactly equivalent representation of all required intermediates.
Arbitrary \emph{selected VB structure spaces} need not be complete.
The two statements are different.

\subsection{Separation of the two spin factors}

Define the same-spin channel actions
\begin{equation}
 D_\sigma=\sum_{PQ}K_{PQ}E_{\sigma P}E_{\sigma Q}.
\end{equation}
They include the particular distribution of one-electron terms implied
by Eq.~\eqref{eq:K}; they are not independently assigned physical
$\alpha$- and $\beta$-spin Hamiltonians. Expanding the two generators gives
\begin{align}
 \Ea_P\Ea_Q(LF\Tr)={}&
 (E_{\alpha P}E_{\alpha Q}L)F\Tr
 +L(E_{\beta P}E_{\beta Q}F)\Tr\notag\\
 &+(E_{\alpha P}L)(E_{\beta Q}F)\Tr
 +(E_{\alpha Q}L)(E_{\beta P}F)\Tr.
\end{align}
The last two terms become identical after summation because $K_{PQ}=K_{QP}$.
Consequently,
\begin{equation}\label{eq:Hfactor}
 \boxed{\ \Ha(LF\Tr)=(D_\alpha L)F\Tr+L(D_\beta F)\Tr
 +2\sum_{PQ}K_{PQ}(E_{\alpha P}L)(E_{\beta Q}F)\Tr.\ }
\end{equation}
The proof uses commutativity of different spin sectors, not commutativity
of generators within one spin sector. It applies to unequal
$U_\alpha,U_\beta$, unequal electron counts, ionic structures, and arbitrary
mixtures of selected structures.

\section{Direct sampling and structure projection}

Introduce the pulled-back factors
\begin{align}\label{eq:pullbacks}
 L_0&=T_\alpha\Tr L,&F_0&=T_\beta\Tr F,\\
 L_H&=T_\alpha\Tr D_\alpha L,&F_H&=T_\beta\Tr D_\beta F,\\
 L_P&=T_\alpha\Tr E_{\alpha P}L,&
 G_P&=T_\beta\Tr\sum_Q K_{PQ}E_{\beta Q}F.
\end{align}
Combining Eqs.~\eqref{eq:T} and \eqref{eq:Hfactor}, the needed original
product entries are
\begin{align}\label{eq:samples}
 m_{ab}&=\sum_{\ell=1}^r(L_0)_{a\ell}(F_0)_{b\ell},\\
 h_{ab}&=\sum_{\ell=1}^r\left[\label{eq:hsamples}
 (L_H)_{a\ell}(F_0)_{b\ell}+(L_0)_{a\ell}(F_H)_{b\ell}
 +2\sum_P(L_P)_{a\ell}(G_P)_{b\ell}\right],
 \qquad (a,b)\in\Omega.
\end{align}
For any Davidson input $\vv v$, build $X=\Wa\vv v$ and return
\begin{equation}\label{eq:final}
 (M\vv v)_I=\sum_{(a,b)\in\Omega}w_{ab,I}m_{ab},\qquad
 (H\vv v)_I=\sum_{(a,b)\in\Omega}w_{ab,I}h_{ab}.
\end{equation}

\begin{proposition}[Exact factored selected-space action]
Equations~\eqref{eq:source}--\eqref{eq:final} evaluate
Eq.~\eqref{eq:restriction} exactly for arbitrary selected structures,
without constructing a complete transformed product array.
\end{proposition}
\begin{proof}
Equation~\eqref{eq:source} is the exact decomposition into nonzero original
columns. Linearity gives Eq.~\eqref{eq:factors}, and expansion of the
Hamiltonian gives Eq.~\eqref{eq:Hfactor}. Applying $\Ta^*$ to each outer
product yields Eqs.~\eqref{eq:samples} and \eqref{eq:hsamples}.
All other final entries have zero
coefficient in $\Wa^*$ by Eq.~\eqref{eq:Wadj}. The final sampling therefore
omits no contribution to Eq.~\eqref{eq:restriction}.
\end{proof}

The contraction can be read as the following sequence:
\begin{equation}
 \vv v\ \xrightarrow{\Wa}\ X=AB\Tr
 \ \xrightarrow{T_\alpha,T_\beta}\ LF\Tr
 \ \xrightarrow{\text{spin generators}}\ \text{factor images}
 \ \xrightarrow{T_\alpha\Tr,T_\beta\Tr}\ (h,m)|_\Omega
 \ \xrightarrow{\Wa^*}\ (H\vv v,M\vv v).
\end{equation}
The middle objects are represented by factors and channel panels, not by
a selected-space approximation to the transformed wave function. This is
how Slater--Condon evaluation remains available without requiring selected
structures to form an invariant space.

\section{Two additional exact reductions}

\subsection{Common spin-exchange parity}

An optional reduction is available when the two one-spin bases and their
orbital maps coincide, and every selected structure has the same certified
transpose parity $\eta=\pm1$. Define
\begin{equation}
 \Pi X=X\Tr,\qquad P_\eta=\tfrac12(\Id+\eta\Pi),\qquad
 \Pi\Wa=\eta\Wa.
\end{equation}
Construct $X_{\rm rep}$ by keeping $a<b$ with coefficient $2X_{ab}$,
keeping diagonal entries with coefficient $X_{aa}$ for $\eta=+1$, and
setting the other entries to zero. Odd parity has zero diagonal. Then
\begin{equation}
 P_\eta X_{\rm rep}=X.
\end{equation}
Since $\Ta$ and a spin-independent $\Ha$ commute with exchange,
their original-basis composite $\mathcal A=\Ta^*\Ha\Ta$ also commutes.
Furthermore $\Wa^*P_\eta=\Wa^*$. It follows that
\begin{equation}
 \Wa^*\mathcal A X_{\rm rep}
 =\Wa^*P_\eta\mathcal A X_{\rm rep}
 =\Wa^*\mathcal A P_\eta X_{\rm rep}
 =\Wa^*\mathcal A X.
\end{equation}
The same proof applies to $\Ta^*\Ta$ for the overlap. Source factorization
may therefore use $X_{\rm rep}$, but the final output sampling still uses
\emph{both} halves of $\Omega$. The unprojected image of $X_{\rm rep}$
is not itself a parity eigenvector.

For a single structure with $m$ independent covalent pairs, this can reduce
the source count from $2^m$ to $2^{m-1}$. It does not imply that the touched
column count always halves for an arbitrary selected space. Equal spin
dimensions alone are insufficient; common parity must be checked from the
structure map. Parity is not a substitute for total-spin adaptation.
Without it, the unreduced factorization remains exact.

\subsection{Causal rows of the exterior adjoint}

Only $A_\Omega$ rows of the $\alpha$ pullbacks and $J$ rows of the
$\beta$ pullbacks enter Eq.~\eqref{eq:samples}. An elementary forward
exterior shear updates $z_j\leftarrow z_j+s\epsilon z_i$, with orbital
coefficient $s$ and fermionic sign $\epsilon=\pm1$. Its adjoint updates
$z_i\leftarrow z_i+s\epsilon z_j$. Starting from the requested final
rows, traverse the adjoint operations in reverse chronological order.
Retain an update exactly when its destination is needed, and then mark
its source as needed. Execute the retained updates in their original
adjoint order, including the diagonal exterior scaling.

For a row selector $Q_A$, this computes $Q_A T_\sigma\Tr Z$ exactly.
Every intermediate row that can influence the answer is retained. The
pruning concerns paths through a linear transformation, not closure of
the selected VB space and not the magnitude of an amplitude. A bitmap on
one-spin rows and retained shear-edge indices are sufficient; no product
array or dense compound matrix is required.

\section{Tiling, storage, and arithmetic}

\subsection{Streaming the exact equations}

Split the $r$ source columns into panels of width $s$, and the $d$ channels
into panels of width $p$. Different source columns are summed independently
in Eq.~\eqref{eq:samples}; the Hamiltonian action never couples source
column labels to one another. For each source panel:
\begin{enumerate}[label=\arabic*.]
 \item Build its sparse $A,B$, and apply $T_\alpha,T_\beta$ to obtain $L,F$.
 \item For a channel panel $\mathcal P$, scatter $E_{\sigma P}$ on the
       corresponding factors.
 \item Stream channel panels $\mathcal Q$ to accumulate
       $V_{\sigma P}=\sum_QK_{PQ}E_{\sigma Q}F_\sigma$,
       where $F_\alpha=L$ and $F_\beta=F$. This is a channel GEMM on
       flattened spin-row/source-column panels.
 \item Gather $\sum_{P\in\mathcal P}E_{\sigma P}V_{\sigma P}$ into
       the same-spin images. Pull back $E_{\alpha P}L$ and $V_{\beta P}$
       and add the cross-spin dot products on $\Omega$.
 \item After all channel panels, pull back the base and same-spin factors,
       add their samples, and reuse the slabs for the next source panel.
\end{enumerate}
Finally apply $\Wa^*$ once to the accumulated support arrays. Source,
channel, and right-hand-side loops reuse shared backing allocations;
threads do not receive private copies of the carrier.

The explicit source/channel numerical slabs require
\begin{equation}\label{eq:memory}
 O\bigl((U_\alpha+U_\beta)sp\bigr)
\end{equation}
storage, plus $O(t)$ support coefficients and images per processed vector.
Additional persistent storage includes the sparse structure map, one-spin
link/exterior metadata, and the $d\times d$ kernel $K$. The source slabs
do not include any $U_\alpha U_\beta$ coefficient array. This is a
workspace statement, not a bound on the total process memory.
For an overlap-only action no channel panels or same-spin Hamiltonian
images are required, and the factor workspace reduces to
$O((U_\alpha+U_\beta)s)$.

\subsection{Where a speed advantage is obtained}

Let $q_\sigma$ count elementary exterior shear edges. With a full channel
panel ($p=d$), a leading-work bound per input vector is
\begin{align}\label{eq:work}
 O\Bigl(&r(U_\alpha+U_\beta)d^2
       +rd(q_\alpha+q_\beta)+trd\notag\\
       &+r\sum_\sigma U_\sigma n_\sigma(n-n_\sigma+1)
       +\operatorname{nnz}(\Wa)\Bigr).
\end{align}
The terms describe channel mixing, exterior pullbacks, final sampling,
Slater--Condon spin links, and the initial/final structure contractions.
The displayed shear count is an upper bound; causal pruning reduces the
selected-row pullbacks. Forward base-factor transforms remain necessary.
For $p<d$, channel mixing keeps the same leading arithmetic order, while
spin scatters can repeat across channel panels by a factor up to
$\lceil d/p\rceil$. A smaller workspace does not change any equation or
variational space.

A generated-product contraction instead has channel-mixing work of order
$N_{\rm cl}d^2$, where $N_{\rm cl}$ counts the generated product
intermediates needed by that contraction order. When this support is broad
and $r\ll\min(U_\alpha,U_\beta)$, the factorized order replaces the
product dimension by a source-weighted sum of one-spin dimensions. For
$U_\alpha\simeq U_\beta=U$ and $N_{\rm cl}\simeq U^2$, the leading
mixing ratio is approximately $U/(2r)$, before other work is counted.

For broad source support, $r$ can approach $U$ and the factorized order
need not be faster. Both orders evaluate the same orthogonal-CI action.
An arithmetic planner may select between them without a memory-based
admission test for the physical selected space. Exact applicability to
arbitrary structures therefore does not mean that their costs are equal
or that selecting one structure universally yields rank one. There is no
claim of eliminating the exponential growth of the complete one-spin
spaces.

\section{Integral adjoints without a product carrier}\label{sec:factor-adjoint}

Let the bra and ket amplitudes have exact factorizations
\begin{equation}
 Y_\ell=L_\ell F_\ell\Tr,\qquad Y_r=L_r F_r\Tr,
 \qquad \mathscr E=\ip{Y_\ell}{\Ha Y_r}.
\end{equation}
Their source counts $r_\ell$ and $r_r$ need not be equal. At fixed factors,
the derivative with respect to the Hamiltonian channel kernel is
\begin{equation}\label{eq:kernel-adjoint}
 \frac{\partial\mathscr E}{\partial K_{PQ}}=
 \ip{\Ea_PY_\ell}{\Ea_QY_r}=G_{PQ}.
\end{equation}
This uses the self-adjoint symmetric generators, but does not assume that
the bra and ket, their source counts, or their spin spaces coincide.

First contract the unacted ket factors into the bra amplitude:
\begin{equation}\label{eq:weighted-bra}
 W_\alpha=Y_\ell F_r=L_\ell(F_\ell\Tr F_r),\qquad
 W_\beta=Y_\ell\Tr L_r=F_\ell(L_\ell\Tr L_r).
\end{equation}
For each channel $P$, form the two channel covectors
\begin{align}\label{eq:channel-covectors}
 C_{\alpha P}&=E_{\alpha P}W_\alpha
            +L_\ell\bigl(F_\ell\Tr(E_{\beta P}F_r)\bigr),\\
 C_{\beta P}&=E_{\beta P}W_\beta
            +F_\ell\bigl(L_\ell\Tr(E_{\alpha P}L_r)\bigr).
\end{align}
Their dimensions are $U_\alpha\times r_r$ and $U_\beta\times r_r$.
Self-adjointness moves the crossed generator from the bra to the ket;
for example, $(E_{\beta P}F_\ell)\Tr F_r
=F_\ell\Tr(E_{\beta P}F_r)$. Thus no channel scatter on the bra
source family is required, even when it contains many panels.
The full kernel derivative is then
\begin{equation}\label{eq:factor-kernel-adjoint}
 \boxed{\quad G_{PQ}=
 \ip{C_{\alpha P}}{E_{\alpha Q}L_r}
 +\ip{C_{\beta P}}{E_{\beta Q}F_r}.\quad}
\end{equation}
To verify this identity, expand $\Ea_QY_r$ as
$(E_{\alpha Q}L_r)F_r\Tr+L_r(E_{\beta Q}F_r)\Tr$ and move each
unacted ket factor to the bra side of the Frobenius product. The resulting
bra contractions are precisely Eq.~\eqref{eq:channel-covectors}. This is
the reverse contraction of the same one-body channel graph used for sigma;
no product-space one-body response or higher-order density tensor is formed.

Only $\operatorname{sym}G=(G+G\Tr)/2$ contributes to symmetric integral
variations. Using the packed diagonal selector $d_P$ and the effective
one-electron matrix $f$ defined above, Eq.~\eqref{eq:K} gives
\begin{equation}
 \delta\mathscr E=
 \tfrac12\ip{\operatorname{sym}G}{\delta V}
 +\sum_P\frac{[(\operatorname{sym}G)d]_P}{N}\,\delta f_P,
 \qquad
 \delta f_{ps}=\delta h_{ps}-\tfrac12\sum_q\delta(pq|qs).
\end{equation}
Unpacking the symmetric pair entries with half weight on each off-diagonal
matrix element gives the full-matrix Frobenius gradients for $h$ and $V$.
The product and factor algorithms share this final kernel pullback in code;
they differ only in the contraction order that generates $G$.

\subsection{A bounded reverse sweep}

For each ket source panel, stream the bra factors to accumulate
$W_\alpha,W_\beta$ once. For each $P$-channel panel, scatter the ket
and these weighted factors, then stream the bra base factors to accumulate
the crossed contributions to $C_{\sigma P}$. Only after the complete bra has been
consumed, stream the ket $Q$ channels and add their Gram matrices to $G$.
Putting the $Q$ contraction inside the bra source loop would repeat the
leading channel GEMM unnecessarily; that order is not used.

Let $s_\ell,s_r$ be the bra and ket panel widths and $p$ the channel width.
The required numerical workspace, measured in doubles, is bounded by
\begin{equation}
 (U_\alpha+U_\beta)
 \bigl[s_\ell+2s_r+(2+\mathbf 1_{p<d})s_rp\bigr]
 +s_\ell s_rp+\max(\omega_\ell s_\ell,\omega_r s_r).
\end{equation}
It includes loaded bra/ket base factors, weighted bra factors, ket channel
panels, channel covectors, and one reusable source moment matrix.
When $p=d$, the ket $P$ panel also supplies all $Q$ channels; otherwise
a separate $Q$ panel is required. Here $\omega_\ell,\omega_r$ count temporary doubles per source
column used by the respective factor loaders, excluding their output views.
They vanish for the original exterior-factor provider; source generation is
sequential, so the larger loading scratch, not their sum, enters the bound.
The immutable metadata and the required returned $d\times d$ gradient are
separate. Both source
and channel widths can decrease to respect the budget, without changing
the represented amplitudes or a derivative tolerance.
Source and channel widths are chosen jointly by the repeated-work estimate
in Section~\ref{sec:factor-schedule}. A complete bra is retained when that
schedule makes it worthwhile; otherwise the same equations stream its panels.

With complete channel panels, the leading dense arithmetic is
\begin{equation}
 O\left((U_\alpha+U_\beta)
    [r_r d^2+r_\ell r_r d]\right),
\end{equation}
instead of product-carrier Gram work of order $U_\alpha U_\beta d^2$.
The source-pair term is essential and is not assumed linear in rank.
Self-adjointness permits exchanging the bra and ket before contraction,
placing the smaller source family on the ket side. One-spin link
traversals and panel regeneration also contribute; a smaller channel
panel can repeat these traversals. No universal speedup is inferred from
the leading dense terms alone.

For an exchange-parity reduction, a representative cannot replace both
amplitudes without projection. A physical bra with parity $\eta$ may,
however, contract an unprojected ket representative $Z$ when
$Y_r=P_\eta Z$ and the integral derivative commutes with exchange:
$\ip{Y_\ell}{\delta\Ha Z}=\ip{Y_\ell}{\delta\Ha Y_r}$.
The caller must establish this identity; the common factor kernel does
not infer a parity projection from equal spin dimensions.

\subsection{Directional factors and the integral part of an HVP}

Let $R$ be the nonsingular upper-triangular orbital map in the exterior
transforms and let $\delta R$ be its direction. Define the spatial generator
and its one-spin representation by
\begin{equation}
 \kappa=\delta R R^{-1},\qquad
 \mathcal G_\sigma(\kappa)=\sum_{pq}\kappa_{pq}
       a^\dagger_{p\sigma}a_{q\sigma}.
\end{equation}
The derivative of an exterior representation obeys
$\delta T_\sigma=\mathcal G_\sigma(\kappa)T_\sigma$.
This generator differentiates the CI representation; it does not replace
the physical nonorthogonal VB orbital coordinates or their finite update.
For $Y=LF\Tr$, with selected structure coefficients allowed to vary,
\begin{equation}\label{eq:factor-tangent}
 \boxed{\quad\delta Y=
   (\mathcal G_\alpha(\kappa)L)F\Tr
   +L(\mathcal G_\beta(\kappa)F)\Tr
   +T_\alpha\,\delta X\,T_\beta\Tr,\qquad
   \delta X=\Wa\,\delta\vv C.\quad}
\end{equation}
The structure map $\Wa$ and its exact support plan are fixed. The last term
therefore uses the same support factorization with the structure-response
coefficients; it is not a projection of either orbital derivative onto the
original support. The two orbital terms generally leave that support and
retain every contribution through their transformed spin factors.

Each term of Eq.~\eqref{eq:factor-tangent} is another factor source.
Concatenating them adds their outer products exactly, without SVD truncation
or a closure assumption. A source panel loads its base factors and applies
the oriented Slater--Condon generator to the required spin factor only.
Unlike a Hamiltonian channel, $\kappa$ need not be symmetric: replacing it by
its symmetric part would give the wrong orbital derivative. The one-factor
temporary uses at most $\max(U_\alpha,U_\beta)s$ doubles for panel width $s$;
it is included in the loader workspace term of the preceding bound.

Write $\mathcal J(Y_\ell,Y_r)$ for the symmetric integral adjoint of
$\langle Y_\ell,\mathcal H Y_r\rangle_{\rm F}$, including the shared
$K$-kernel pullback. Since each physical integral variation is self-adjoint,
\begin{equation}\label{eq:factor-adjoint-tangent}
 \delta\mathcal J(Y,Y)=2\mathcal J(\delta Y,Y).
\end{equation}
Thus the integral-gradient derivative needed in an orbital HVP can be
contracted from the directional factor source directly, without constructing
$\delta Y$ or any higher-order RDM for this contraction. If the accepted
and directional factor counts are $r$ and $r_\delta$, the leading dense work is
\begin{equation}
 O\left((U_\alpha+U_\beta)
  [\min(r,r_\delta)d^2+r\,r_\delta d]\right),
\end{equation}
plus the one-spin link traversals and source transforms. The count
$r_\delta$ includes both orbital terms and any structure response, including
the physical parity reconstruction when required. It must not be estimated
from the accepted source count alone.

Equation~\eqref{eq:factor-adjoint-tangent} is the integral part of the HVP,
not its complete orbital pullback. The derivative of the exterior-map
adjoint, the full carrier residual, and the state-average gauge terms remain
required. The present implementation keeps those terms while migrating the
integral contraction to Eq.~\eqref{eq:factor-adjoint-tangent}.

\subsection{Oriented generator and residual contractions}\label{sec:factor-generator}

Define the orbital matrix $\mathcal D(X,Y)$ by
\begin{equation}
 [\mathcal D(X,Y)]_{pq}=
 \langle X,(a^\dagger_{p\alpha}a_{q\alpha}
                 +a^\dagger_{p\beta}a_{q\beta})Y\rangle_{\rm F}.
\end{equation}
It is the coefficient of $\kappa_{pq}$ in
$\langle X,\mathcal G(\kappa)Y\rangle_{\rm F}$ and need not be symmetric.
For $Y=L_rF_r\Tr$, form only the weighted single-spin panels
\begin{equation}
 W_\alpha=XF_r,\qquad W_\beta=X\Tr L_r.
\end{equation}
Then
\begin{equation}\label{eq:generator-factor}
 [\mathcal D(X,Y)]_{pq}=
 \sum_j\langle (W_\alpha)_j,
      a^\dagger_{p\alpha}a_{q\alpha}(L_r)_j\rangle
 +\sum_j\langle (W_\beta)_j,
      a^\dagger_{p\beta}a_{q\beta}(F_r)_j\rangle.
\end{equation}
If $X=L_\ell F_\ell\Tr$, each weighted panel is obtained from source-pair
moments, $XF_r=L_\ell(F_\ell\Tr F_r)$ and
$X\Tr L_r=F_\ell(L_\ell\Tr L_r)$. The final contractions traverse the
oriented single-spin occupation/replacement links. They are not symmetric
Hamiltonian-channel derivatives. For real coefficients,
$\mathcal D(X,Y)=\mathcal D(Y,X)\Tr$; this permits a transposed
contraction order but not an unqualified interchange of sources.

For the residual $Z=(\Ha-E)X$, the required panels can also be obtained
without $Z$ or $\Ha X$. Using the same-spin actions $D_\sigma$ and
the ket channel mixtures
\begin{equation}
 V_{\alpha P}=\sum_QK_{PQ}E_{\alpha Q}L_r,\qquad
 V_{\beta P}=\sum_QK_{PQ}E_{\beta Q}F_r,
\end{equation}
the exact residual-weighted panels are
\begin{align}\label{eq:residual-panels}
 Z_\alpha={}&D_\alpha W_\alpha+X(D_\beta F_r)
       +2\sum_P E_{\alpha P}(X V_{\beta P})-E W_\alpha,\\
 Z_\beta={}&D_\beta W_\beta+X\Tr(D_\alpha L_r)
       +2\sum_P E_{\beta P}(X\Tr V_{\alpha P})-E W_\beta.
\end{align}
Thus $Z_\alpha=ZF_r$ and $Z_\beta=Z\Tr L_r$. Replacing
$W_\alpha,W_\beta$ by $Z_\alpha,Z_\beta$ in
Eq.~\eqref{eq:generator-factor} gives $\mathcal D((\Ha-E)X,Y)$.
The spectator actions have been moved to the ket using self-adjointness
of $D_\sigma$; the crossed term follows from Eq.~\eqref{eq:Hfactor}.
Self-adjointness of $\Ha$ does \emph{not} justify exchanging the two
sources here, because $\Ha$ generally does not commute with an oriented
one-body generator.

The implementation reuses the factor/channel sweep twice: once for the
weighted panels to produce $D_\sigma W_\sigma$, and once for the ket
panels to produce the spectator and crossed terms. After same-spin gather,
the ket's raw channel buffers are reused for $X V_{\beta P}$ and
$X\Tr V_{\alpha P}$, then consumed by the crossed gather. No separate
channel scatter/mixing implementation or product-space residual is required.

With complete channel panels, the leading dense arithmetic is
$O((U_\alpha+U_\beta)[r_r d^2+r_\ell r_r d])$, plus the oriented
single-spin link contractions. Without $\Ha$, it reduces to
\begin{equation}
 O\left((U_\alpha+U_\beta)r_\ell r_r
  +r_r\sum_{\sigma}U_\sigma n_\sigma(n-n_\sigma+1)\right).
\end{equation}
Accepted-state transition densities use this estimate to choose between
factor moments and product-space links. Their planner does not reuse the
Hamiltonian-channel cost estimate.

A conservative numerical workspace bound for the residual contraction,
in doubles, is
\begin{equation}
 (U_\alpha+U_\beta)[s_\ell+(3p+9)s_r]
 +s_\ell s_r\max(1,p)
 +\max(\omega_\ell s_\ell,\omega_r s_r)+(t+1)n^2,
\end{equation}
where $t$ bounds the number of workers and the last term covers the small
orbital-matrix reduction and its temporary result. The required returned
$n\times n$ matrix and immutable metadata are separate. For a generator
without $\Ha$, the channel sweep and residual outputs are absent, giving
$(U_\alpha+U_\beta)(s_\ell+2s_r)+s_\ell s_r$ plus the same loading
and reduction terms. No per-worker spin/product array is allocated.

\subsection{Joint source/channel scheduling}\label{sec:factor-schedule}

For source widths $s_\ell,s_r$ and channel width $p$, define
\begin{equation}
 N_\ell=\left\lceil r_\ell/s_\ell\right\rceil,\qquad
 N_r=\left\lceil r_r/s_r\right\rceil,\qquad
 N_P=\left\lceil d/p\right\rceil,\qquad U=U_\alpha+U_\beta.
\end{equation}
The leading dense channel and moment arithmetic does not depend on these
widths, but source regeneration, link traversal and output traffic do.
In particular, the number of loaded bra source columns is
\begin{equation}
 C_\ell=\begin{cases}
 r_\ell,&s_\ell=r_\ell,\\
 N_r(N_P+b)r_\ell,&s_\ell<r_\ell,
 \end{cases}\qquad
 b=\begin{cases}1,&\text{integral adjoint},\\2,&\text{residual generator}.
 \end{cases}
\end{equation}
Here $b$ counts the channel-independent moment sweeps. If generating a
source column costs $\ell_\ell$ nominal scalar operations and element
touches, source work is $C_\ell\ell_\ell$. Each provider declares this work
from its sparse writes, diagonal scales and exterior shear multiply--adds;
directional providers add the oriented generator work. This metadata affects
only the evaluation schedule, never the amplitudes, ranks or tolerances.

The single-replacement links are generated in ascending packed-channel
order. For each larger orbital index, ascending lower indices with opposite
occupation generate every link once. A channel tile then uses binary ranges
in the existing link and occupation arrays. It does not scan all links and
reject those outside the tile, nor allocate a string-by-channel offset table.
Let $J=\sum_\sigma U_\sigma n_\sigma(n-n_\sigma+1)$ count the single-spin
occupation/replacement links. Across all source columns, matching-link visits
scale as $r_rv(N_P)J$, where
\begin{equation}
 v_I(N_P)=N_P+1,\qquad v_R(N_P)=2N_P+3.
\end{equation}
Without channel ranges these counts were respectively $N_P^2+N_P$ and
$2N_P^2+3N_P$. Binary lookup comparisons still repeat with those latter
counts and are included through an upper comparison estimate $B$ per
complete single-spin row family. The complete-channel case needs no binary
lookup penalty.

The schedule also counts repeated output reads and writes in the moment
GEMMs, and repeated channel zero fills:
\begin{align}
 W_{\rm out}&=2U r_r(d+b)N_\ell,\\
 W_{0,I}&=U r_r d(N_P+1),&
 W_{0,R}&=2U r_r dN_P.
\end{align}
Combining these terms gives the nominal variable-work estimate
\begin{equation}
 \widehat W=C_\ell\ell_\ell+
   r_r\bigl[3v(N_P)J+s(N_P)B\bigr]+W_0+W_{\rm out},
\end{equation}
where $s_I=N_P^2+N_P$, $s_R=2N_P^2+3N_P$, and the factor three counts
one matching-link visit plus its multiply and add. Terms invariant across
the candidate schedules are omitted. The estimate combines scalar work and
element traffic without fitted hardware weights; it is not a prediction of
absolute wall time, which requires native measurement.

The planner searches feasible source and channel widths under the same memory
bound. For each ket/channel choice it takes the largest feasible bra panel;
ties use fewer moment panels. Partial-channel residual sweeps currently use
one ket column because the shared forward workspace and moment buffers must
use the same channel width. Full-channel sweeps allow wider ket panels.
This is a workspace scheduling limit, not a restriction on the selected VB space.

\section{Correspondence with the current implementation}

The original column contraction corresponds to revision \texttt{a5004add};
Section~\ref{sec:cover} adds the minimum mixed-support decomposition.
Paths below are relative to \texttt{src/}; this table records the
implementation boundary rather than introducing physical assumptions.
\begin{center}\small
\begin{tabular}{>{\raggedright\arraybackslash}p{0.37\linewidth}
                >{\raggedright\arraybackslash}p{0.54\linewidth}}
\toprule
Mathematical operation & Implementation \\
\midrule
Original support and minimum source cover &
\path{vbscf/structures/assembly/sources.cpp} \\
VB source generation, sampling and $\Wa^*$ &
\path{vbscf/structures/assembly/source_tiles.cpp}\newline
\texttt{apply\_source\_tiles()} \\
Common factor sweep and integral adjoint &
\path{mcscf/ci/factors.cpp} \\
Generator and residual-weighted factor adjoints &
\path{mcscf/ci/factor_generator.cpp} \\
Joint source/channel work plan &
\path{mcscf/ci/factor_plan.cpp} \\
Hamiltonian channel kernel, Eq.~\eqref{eq:K} &
\path{mcscf/ci/sigma.cpp}\newline
\texttt{build\_hamiltonian\_pair\_kernel()} \\
Slater--Condon spin connections &
\path{mcscf/ci/spin_channels.cpp} \\
Exterior transforms and causal adjoint rows &
\path{vbscf/structures/orthogonal_ci/exterior_transform.cpp}\newline
\texttt{projection()} and \texttt{apply\_adjoint\_left\_block()} \\
Independent complete-carrier comparisons &
\path{../tests/vbscf/unit/test_structure_tiles.cpp} \\
\bottomrule
\end{tabular}
\end{center}

Fixed-point structure Hamiltonian/overlap actions use the common factor
sweep. Accepted-state integral gradients choose between the exact factor and
product contractions by their arithmetic costs. Raw state coefficients are
retained in structure space. The orthogonal state, its full carrier residual,
and generator-response consumers still require migration. Directional
integral adjoints use Eq.~\eqref{eq:factor-adjoint-tangent} when the factor
order has less estimated arithmetic. A selected-space eigenvector need not be
stationary on the full carrier: its residual contribution to the
exterior-transform gradient is retained.
Accepted-state transition densities use Eq.~\eqref{eq:generator-factor}
when advantageous. The residual-panel kernel of
Eq.~\eqref{eq:residual-panels} is available and independently tested;
accepted-state residual storage is not yet removed from the remaining consumers.

\paragraph{Main result.}
Equations~\eqref{eq:Hfactor}--\eqref{eq:final} give the exact
selected-space restriction through spin factors and Slater--Condon actions,
without a closure assumption or a complete transformed product array.

\end{document}
